\chapter{Arsitektur Modern dan Model Generatif}
\label{ch:arsitektur}

Kuliah lanjutan deep learning saya ikuti pada musim panas 2025, dan latihannya setahun kemudian.
Tugas-tugas latihannya membentuk sebuah tangga: PyTorch dan training loop, jaringan visi dan
transfer learning, pemantauan eksperimen dan CNN lanjutan, autoencoder dan VAE, contoh adversarial
dan GAN, lalu normalizing flow. Bab ini menaiki tangga yang sama.

\Cref{ch:dlbasic} membangun jaringan dari bata-bata dasarnya. Bab ini membahas dua hal yang
membuat deep learning modern berbeda: arsitektur yang memungkinkan jaringan amat dalam dilatih
dengan stabil, dan model yang tidak hanya mengenali data tetapi juga membangkitkan data baru.

\section{Dari LeNet ke ResNet}

LeNet \citep{lecun1998} dan AlexNet \citep{krizhevsky2017} masih memakai kernel yang relatif besar.
Keluarga VGG hanya memakai kernel $3\times3$ yang ditumpuk. Dua lapisan $3\times3$ melihat daerah
$5\times5$, tiga lapisan melihat $7\times7$, dengan parameter yang lebih sedikit dan nonlinearitas
yang lebih banyak. Dengan 64 kanal masuk dan keluar, satu lapisan $5\times5$ membutuhkan
$25\cdot64^2=102\,400$ bobot, sedangkan dua lapisan $3\times3$ dengan receptive field yang sama
cukup dengan $2\cdot9\cdot64^2=73\,728$ bobot, sekitar 28 persen lebih sedikit, dan memberi dua
nonlinearitas alih-alih satu.

Menambah lapisan seharusnya tidak merugikan, karena lapisan tambahan bisa saja belajar menjadi
fungsi identitas. Kenyataannya, jaringan polos yang dalam justru punya galat \emph{latih}
lebih tinggi, jadi masalahnya terletak pada optimasi, bukan overfitting. ResNet \citep{he2016}
mengubah tugas setiap blok dari mempelajari $H(\vx)$ menjadi mempelajari selisihnya,
\begin{equation}
\vy=\vx+F(\vx).
\end{equation}
Kalau identitas sudah cukup baik, blok cukup mendorong $F$ ke nol. Gradien pun mendapat jalan tol,
$\partial\vy/\partial\vx=\bm I+\partial F/\partial\vx$, sehingga sinyal mundur tidak harus melewati
semua nonlinearitas. Gagasannya sama dengan constant error carousel pada LSTM (\cref{ch:lstm}),
hanya saja kali ini di sepanjang kedalaman, bukan di sepanjang waktu.

Editor yang baik tidak menulis ulang seluruh esai setiap kali merevisi. Ia menulis catatan
perubahan di pinggir halaman, dan kalau esainya sudah bagus, catatannya kosong. Blok residual
bekerja seperti editor itu: ia hanya menambahkan koreksi pada apa yang sudah ada.

Latihan membahas dua variasi. Pada ResNet \istilah{pre-activation}, normalisasi dan aktivasi
diletakkan sebelum konvolusi, sehingga bloknya berbentuk $\vx+F(\phi(\vx))$ alih-alih
$\phi(\vx+F(\vx))$. Ketika jalur pintas dan jalur utama berbeda jumlah kanal, konvolusi $1\times1$
dipakai untuk menyamakannya.

Ada juga blok yang dirancang untuk efisiensi. \istilah{Depthwise-separable convolution} memecah
konvolusi menjadi dua langkah. Setiap kanal dikonvolusi sendiri-sendiri, lalu kanal-kanal dicampur
dengan konvolusi $1\times1$. Untuk kernel $3\times3$ dan 64 kanal, jumlah bobot turun dari
$9\cdot64\cdot64=36\,864$ menjadi $9\cdot64+64\cdot64=4\,672$. Blok \istilah{squeeze-and-excitation}
merata-ratakan setiap kanal menjadi satu angka, melewatkannya ke MLP kecil, lalu memakai hasilnya
untuk menskalakan kanal, sehingga jaringan belajar kanal mana yang penting untuk gambar tertentu.
\istilah{Vision transformer} memotong gambar menjadi petak-petak $P\times P$ dan memperlakukan setiap
petak sebagai token \citep{dosovitskiy2021}. Gambar $224\times224$ dengan $P=16$ menghasilkan
$14\times14=196$ token. MLP-Mixer bergantian mencampur informasi antar token dan antar kanal
dengan MLP biasa, tanpa konvolusi dan tanpa attention.

Melatih jaringan besar dari nol membutuhkan data besar. \istilah{Transfer learning} memakai jaringan
yang sudah dilatih pada ImageNet, membekukan sebagian besar lapisannya, lalu melatih ulang
kepala klasifikasinya saja. Fitur umum seperti tepi dan tekstur sudah dipelajari, sehingga data
yang sedikit pun cukup untuk tugas baru.

\section{Dari mengenali ke membangkitkan}

Model diskriminatif mempelajari $p(y\mid\vx)$. Model generatif mempelajari $p(\vx)$ itu sendiri,
sehingga bisa membangkitkan contoh baru, mengisi bagian yang hilang, atau menilai seberapa wajar
sebuah data. Kuliah membahas lima keluarga.

\subsection{Autoencoder}

Autoencoder memampatkan $\vx$ menjadi kode $\vz=e(\vx)$ yang berdimensi kecil, lalu
merekonstruksinya menjadi $\hat\vx=d(\vz)$. Loss rekonstruksinya mengikuti asumsi noise
(\cref{tab:noiseloss}): MSE untuk data kontinu yang dianggap Gauss, binary cross-entropy dengan
keluaran sigmoid untuk data di antara nol dan satu. Autoencoder linear dengan loss kuadrat
menemukan subruang yang sama dengan PCA. Kelemahannya, autoencoder biasa tidak menjamin ruang kode
yang rapi. Titik acak di ruang kode bisa didekode menjadi sesuatu yang tidak bermakna.

\subsection{Variational autoencoder}

VAE \citep{kingma2014} memperbaiki hal itu dengan memberi struktur probabilistik. Encoder
mengeluarkan distribusi $q_\phi(\vz\mid\vx)=\N(\bm\mu,\mathrm{diag}\,\bm\sigma^2)$, decoder memodelkan
$p_\theta(\vx\mid\vz)$, dan ruang kode diberi prior $p(\vz)=\N(\bm0,\bm I)$. Pelatihannya
memaksimalkan \istilah{evidence lower bound}:
\begin{equation}
\log p_\theta(\vx)\ \ge\ \E_{q_\phi(\vz\mid\vx)}\big[\log p_\theta(\vx\mid\vz)\big]
-\KL\big(q_\phi(\vz\mid\vx)\,\|\,p(\vz)\big).
\label{eq:elbo6}
\end{equation}
Suku pertama mengukur rekonstruksi, suku kedua mengukur kedekatan ke prior. Untuk Gauss diagonal,
divergensi KL punya bentuk tertutup,
$\KL=\tfrac12\sum_j\big(\sigma_j^2+\mu_j^2-1-\log\sigma_j^2\big)$. Untuk satu dimensi, rumus ini
didapat dari definisi $\KL(q\|p)=\E_q[\log q-\log p]$ dengan $q=\N(\mu,\sigma^2)$ dan $p=\N(0,1)$. Log
kerapatan keduanya adalah $\log q=-\tfrac12\log(2\pi\sigma^2)-(z-\mu)^2/(2\sigma^2)$ dan
$\log p=-\tfrac12\log2\pi-z^2/2$. Di bawah $q$, $\E[(z-\mu)^2]=\sigma^2$ dan $\E[z^2]=\sigma^2+\mu^2$, sehingga
$\E_q[\log q-\log p]=-\tfrac12\log\sigma^2-\tfrac12+\tfrac12(\sigma^2+\mu^2)$, yang sama dengan rumus di
atas. Untuk dimensi yang independen, KL-nya tinggal dijumlahkan. Satu dimensi laten dengan
$\mu=1$ dan $\sigma=1$ memberi KL sebesar $0{,}5$. Dengan $\mu=0$ dan $\sigma=0{,}5$, hasilnya
$\tfrac12(0{,}25-1+1{,}386)\approx0{,}318$. KL bernilai nol hanya kalau $\mu=0$ dan $\sigma=1$, yaitu
tepat sama dengan prior.

Untuk membawa gradien melewati pengambilan sampel acak, VAE memakai \istilah{reparameterization trick}.
Sampel ditulis sebagai $\vz=\bm\mu+\bm\sigma\odot\bm\epsilon$ dengan $\bm\epsilon\sim\N(\bm0,\bm I)$.
Keacakan dipindahkan ke $\bm\epsilon$, sehingga $\vz$ menjadi fungsi yang bisa diturunkan terhadap
$\bm\mu$ dan $\bm\sigma$.

Saya membayangkan encoder sebagai orang yang melihat foto lalu menulis deskripsi singkat yang
sengaja agak samar, misalnya ``wajah bulat, rambut pendek, kira-kira tersenyum''. Decoder adalah
pelukis yang menggambar dari deskripsi itu. Suku KL memaksa semua deskripsi memakai kosakata yang
sama dan tidak terlalu spesifik, sehingga deskripsi baru yang dikarang bebas pun tetap
menghasilkan wajah yang masuk akal.

Latihan menekankan satu kenop praktis: bobot keseimbangan antara rekonstruksi dan KL. Bobot KL yang
terlalu besar membuat gambar kabur dan ruang laten tidak terpakai. Bobot yang terlalu kecil
membuat VAE kembali menjadi autoencoder biasa. Penurunan ELBO yang lebih teliti, termasuk versi
hierarkis yang mengarah ke model difusi, ada di \cref{ch:mlat}.

\subsection{Generative adversarial networks}

GAN \citep{goodfellow2014} mempertandingkan dua jaringan. Generator $G$ mengubah noise $\vz$ menjadi
sampel palsu, dan diskriminator $D$ menebak apakah sebuah sampel asli atau palsu:
\begin{equation}
\min_G\max_D\ \E_{\vx\sim p_{\mathrm{data}}}\big[\log D(\vx)\big]+\E_{\vz}\big[\log\big(1-D(G(\vz))\big)\big].
\end{equation}
Bentuk ini tidak lain adalah binary cross-entropy diskriminator yang tandanya dibalik untuk
generator. Di awal pelatihan, diskriminator mudah menang dan $\log(1-D(G(\vz)))$ jenuh, sehingga
gradien generator kecil. Jalan keluarnya adalah \istilah{non-saturating loss}, di mana generator
memaksimalkan $\log D(G(\vz))$ saja.

Apa yang sebenarnya dioptimalkan permainan ini? Untuk generator yang tetap, dengan distribusi sampel
palsu $p_g$, tujuan diskriminator bisa ditulis sebagai integral
$\int\big[p_{\mathrm{data}}(\vx)\log D(\vx)+p_g(\vx)\log(1-D(\vx))\big]\dd\vx$. Di setiap titik $\vx$,
fungsi $a\log D+b\log(1-D)$ maksimal pada $D=a/(a+b)$, sehingga diskriminator terbaik adalah
\begin{equation}
D^*(\vx)=\frac{p_{\mathrm{data}}(\vx)}{p_{\mathrm{data}}(\vx)+p_g(\vx)}.
\end{equation}
Memasukkan $D^*$ kembali ke tujuan permainan memberi $-\log4+2\,\mathrm{JS}(p_{\mathrm{data}}\|p_g)$, dengan
JS divergensi Jensen--Shannon \citep{goodfellow2014}. Jadi generator, ketika melawan diskriminator yang
sempurna, sebenarnya meminimalkan jarak antara distribusi data dan distribusi sampelnya, dan
minimumnya tercapai tepat ketika keduanya sama, saat $D^*=\tfrac12$ di mana-mana: diskriminator hanya bisa
menebak dengan melempar koin.

Pemalsu uang belajar mencetak uang yang lolos pemeriksaan, sementara detektif belajar membedakan
uang asli dan palsu. Setiap kali detektif menemukan cacat baru, pemalsu memperbaikinya. Kalau
keduanya seimbang, uang palsu akhirnya tidak bisa dibedakan dari yang asli. Masalah khas GAN
adalah \istilah{mode collapse}, ketika generator menemukan beberapa sampel yang selalu lolos lalu
berhenti menghasilkan variasi lain. Pelatihannya juga tidak stabil, karena tidak ada satu fungsi
loss yang turun secara monoton.

\subsection{Normalizing flows}

Flow membangun distribusi yang rumit dengan mengubah distribusi sederhana melalui fungsi yang bisa
dibalik \citep{rezende2015}. Kalau $\vx=g(\vz)$ dengan $g$ bijektif, rumus perubahan variabel
memberi kerapatan yang eksak,
\begin{equation}
\log p_X(\vx)=\log p_Z\big(g^{-1}(\vx)\big)+\log\Big|\det\frac{\partial g^{-1}}{\partial\vx}\Big|.
\end{equation}
Dalam satu dimensi, rumus ini berasal dari kekekalan peluang. Peluang yang jatuh di selang kecil
$[z,z+\dd z]$ harus sama dengan peluang di selang padanannya $[x,x+\dd x]$, sehingga
$p_X(x)|\dd x|=p_Z(z)|\dd z|$, atau $p_X(x)=p_Z(z)\,|\dd z/\dd x|$. Kalau $x=2z$ dengan $z\sim\N(0,1)$, setiap
selang diregangkan dua kali, sehingga kerapatannya harus turun separuh: $p_X(x)=\tfrac12p_Z(x/2)$, yang
memang kerapatan $\N(0,4)$. Determinan Jacobian adalah versi banyak dimensi dari faktor regangan ini,
dan ia mengukur perubahan volume lokal. Untuk dimensi besar, determinan umum mahal
dihitung, jadi lapisan flow dirancang agar Jacobian-nya segitiga. \istilah{Coupling layer} membagi
vektor menjadi dua bagian, membiarkan bagian pertama apa adanya, dan mengubah bagian kedua dengan
fungsi dari bagian pertama. Untuk coupling aditif, $\vx_2'=\vx_2+t(\vx_1)$, determinannya bernilai
satu dan inversnya analitis. Syarat bijeksi berarti dimensi tidak boleh berubah. Imbalannya,
flow memberi likelihood yang eksak, bukan batas bawah seperti ELBO.

\subsection{Model difusi dan model energi}

Model difusi \citep{sohldickstein2015,ho2020} merusak data perlahan-lahan dengan noise Gauss, lalu
melatih jaringan untuk membalik proses itu langkah demi langkah. Mereka kini menjadi model
generatif gambar yang paling kuat, dan matematikanya, yang melibatkan persamaan diferensial
stokastik dan persamaan Fokker--Planck, dibahas lengkap di \cref{ch:mlat}. Model berbasis energi,
termasuk mesin Boltzmann, mendefinisikan $p(\vx)\propto\exp(-E(\vx))$. Mudah ditulis, tetapi sulit
dilatih, karena konstanta normalisasinya adalah integral atas seluruh ruang.

\Cref{tab:generatif} merangkum kelima keluarga ini.

\begin{table}[ht]
\centering\small
\caption{Lima keluarga model generatif.}
\label{tab:generatif}
\begin{tabularx}{\linewidth}{@{}lXXX@{}}
\toprule
Model & Dilatih dengan & Kekuatan & Kelemahan \\
\midrule
VAE & ELBO & ruang laten terstruktur, stabil & sampel cenderung kabur \\
GAN & permainan minimax & sampel tajam & tidak stabil, mode collapse \\
Flow & likelihood eksak & kerapatan eksak, bisa dibalik & arsitektur terbatas, dimensi tetap \\
Difusi & denoising & kualitas dan keragaman & sampling lambat \\
Energi & beragam, sering MCMC & fleksibel & normalisasi sulit \\
\bottomrule
\end{tabularx}
\end{table}

\section{Contoh adversarial}

Jaringan yang akurat bisa dibohongi dengan perubahan yang tidak terlihat oleh mata.
\istilah{Fast gradient sign method} menggeser input ke arah yang paling menaikkan loss,
\begin{equation}
\vx'=\vx+\epsilon\,\mathrm{sign}\big(\nabla_{\vx}\Loss(\vx,y)\big).
\end{equation}
\istilah{Projected gradient ascent} mengulang langkah ini beberapa kali, dan setiap kali
memproyeksikan hasilnya kembali ke bola $\|\vx'-\vx\|_\infty\le\epsilon$ serta ke rentang piksel
yang sah. Serangan seperti ini bekerja juga di dunia fisik \citep{eykholt2018}. Melatih model dengan
contoh adversarial membuatnya lebih tahan, dengan harga akurasi yang sedikit turun pada data
bersih.

\section{Sedikit teori}

Naskah kuliah ditutup dengan beberapa hasil teoretis yang menjelaskan mengapa semua ini bekerja.
Jaringan dalam paling mudah dilatih kalau nilai singular Jacobian input--output mendekati satu,
sehingga sinyal maju dan gradien mundur sama-sama tidak menyusut atau meledak. Sifat ini disebut
isometri dinamis, dan ia adalah versi matriks dari argumen varians di \cref{ch:dlbasic}. Pada
lebar tak hingga dan laju belajar kecil, pelatihan jaringan berperilaku seperti regresi kernel
dengan kernel yang tetap, yaitu \istilah{neural tangent kernel} \citep{jacot2018}. Teori ini
menjelaskan konvergensi jaringan yang lebar sekali dan, seperti akan kita lihat di \cref{ch:pinn},
juga kegagalan pelatihan PINN. Sebelum dilatih, keluaran jaringan acak yang lebar mendekati
proses Gauss, sebuah hubungan yang menjembatani deep learning dengan statistik Bayesian
\citep{rasmussen2006}.

Jaringan biasa hanya memberi satu tebakan, tanpa keterangan seberapa yakin ia. MC dropout
menjalankan dropout saat inferensi dan memperlakukan variasi keluarannya sebagai ketidakpastian
\citep{gal2016}. Deep ensembles melatih beberapa jaringan dari inisialisasi yang berbeda dan melihat
seberapa jauh mereka berbeda pendapat \citep{lakshminarayanan2017}.

\section{Di luar kelas}

Arsitektur dan model generatif di bab ini menyebar ke hampir semua bidang. U-Net lahir untuk citra
biomedis dan tetap menjadi standar segmentasi organ dan tumor (\cref{ch:medis}). Transfer learning
membuat jaringan visi yang dilatih pada foto sehari-hari berguna untuk citra mikroskop dan citra medis
yang datanya sedikit. VAE dan model generatif berbasis teks dipakai untuk merancang molekul baru, dan
metrik seperti Fréchet ChemNet Distance meminjam cara mengevaluasi GAN untuk kimia (\cref{ch:hayati}).
Contoh adversarial mengingatkan bahwa sensor dan sistem persepsi kendaraan otonom bisa diserang
(\cref{ch:robot}).

Di simulasi, untuk medan pada grid teratur, U-Net adalah pilihan bawaan sebagai model pengganti. Ia
menggabungkan konteks yang luas, penting karena tekanan di satu titik dipengaruhi seluruh domain,
dengan detail lokal seperti lapisan batas. Blok residual cocok untuk model langkah waktu, karena
$\vu_{t+1}=\vu_t+F(\vu_t)$ tidak lain adalah integrator Euler dengan $F$ yang dipelajari.

Encoder sebuah autoencoder bisa memetakan medan aliran ke ruang laten berdimensi kecil.
Dinamikanya lalu dipelajari di ruang laten itu, dan decoder mengembalikannya ke medan penuh. Ini
salah satu bentuk model tereduksi yang dibahas di \cref{ch:penemuan}.

Turbulensi bersifat chaos, sehingga lebih jujur memprediksi distribusi keadaan masa depan daripada
satu keadaan tunggal. Model difusi kini dipakai untuk prakiraan cuaca probabilistik
\citep{price2025} dan untuk simulasi aliran turbulen \citep{kohl2023}. Model pengganti yang dipakai
untuk keputusan teknik juga harus bisa mengatakan bahwa ia tidak yakin, dan ensemble serta model
generatif memberi kemampuan itu.

\sumberbab{Bab ini mengikuti bagian kedua naskah kuliah Deep Learning and Neural Nets II, yang kini
bernama Deep Learning: Architectures and Generative Techniques (arsitektur lanjut, model generatif,
teori deep learning, optimasi orde tinggi, Bayesian deep learning), beserta enam tugas latihannya.
Bacaan pendamping: \citet{goodfellow2016}, \citet{prince2023}, dan \citet{kingma2019}.}
